\documentclass[%
 reprint,
 amsmath,amssymb,
 aps,
]{revtex4-1}

\usepackage{graphicx}% Include figure files
\usepackage{dcolumn}% Align table columns on decimal point
\usepackage{bm}% bold math
\usepackage[dvipdfm,colorlinks=true,linkcolor=blue,filecolor=blue, urlcolor=blue]{hyperref}
\usepackage{braket}
\usepackage{caption} % for captionsetup[figure]
\usepackage{threeparttable} % for table caption modification
\usepackage{breqn} % dmath: automatic equation line breaking
\numberwithin{equation}{section}% (p.37)

% \sloppy sets all text to loose spacing, allowing overflow and uneven word spacing;
% not recommended for academic papers requiring strict typesetting.
\tolerance=1000
\emergencystretch=1em
\hyphenpenalty=3000
\exhyphenpenalty=3000
\flushbottom
\usepackage[final]{microtype}
\usepackage{ragged2e}  % Add this package for \justifying command
\usepackage{booktabs}

\begin{document}

    \captionsetup[figure]{justification=raggedright,
        labelfont={bf},name={Fig.},labelsep=period} % raggedright for left-aligned figure captions
    \captionsetup[table]{
        labelfont={bf},name={Table.},labelsep=period}


\title{Anomalous Magnetic Moments Without Fields:\\ A Geometric and Observer-Dependent Interpretation}
\author{Satoshi Hanamura}
\email{hana.tensor@gmail.com}
\date{June 29, 2025}



\begin{abstract}
Building upon the previously established 0-Sphere electron model with its internal structure of two thermal kernels and photon sphere exhibiting Zitterbewegung oscillations, we present a new observer-dependent interpretation that fundamentally reframes the origin of anomalous magnetic moments in quantum electrodynamics. While conventional QED requires external magnetic field interactions for any observable anomaly to manifest, our reinterpretation of the existing 0-Sphere framework demonstrates that anomalous magnetic moments can arise as purely geometric effects in free electrons through relativistic observer transformations. The crucial new insight emerges from recognizing that the same internal Zitterbewegung motion at $v \sim 0.04c$ yields different measurements depending on the observer's reference frame: an observer co-moving with the internal motion measures a $g$-factor of exactly 2 consistent with Dirac theory, while an observer in the laboratory frame perceives Lorentz contraction effects that manifest as the anomalous magnetic moment through the predicted relationship $\gamma = 1 + a_e$. This observer-dependent reinterpretation transforms our understanding of quantum magnetic anomalies from interaction-based corrections to relativistic geometric consequences of internal motion, suggesting that quantum phenomena traditionally attributed to virtual particle interactions may represent observable effects of relativistic internal structure, thereby opening new avenues for understanding quantum mechanics through geometric principles without requiring external field perturbations.
\end{abstract}


\maketitle

% ========================================
% Section 1: Introduction
% ========================================
\section{Introduction}

The anomalous magnetic moment of elementary particles represents one of the most precisely measured quantities in quantum electrodynamics~\cite{fan2023,hanneke2008,gabrielse2006}, yet its fundamental origin remains tied to complex perturbative calculations involving virtual particle interactions~\cite{schwinger1948,feynman1948,feynman1949,tomonaga1946,aoyama2015}. The conventional QED framework interprets the deviation of the electron's $g$-factor from the Dirac prediction~\cite{dirac1928} of 2 as arising exclusively through interactions with external electromagnetic fields, requiring $|\mathbf{B}| > 0$ for any observable anomaly to manifest. This interpretation fundamentally assumes that free electrons in the absence of external fields should exhibit no deviation from the classical Dirac value, treating the anomalous magnetic moment as an interaction-dependent phenomenon rather than an intrinsic property of the electron's internal structure.

Recent developments in the 0-Sphere electron model suggest an alternative interpretation where the anomalous magnetic moment emerges as a geometric consequence of internal relativistic motion observable in truly free electrons~\cite{hanamura2018}. This framework proposes that electrons possess an internal structure consisting of two thermal kernels exchanging energy through a photon sphere exhibiting Zitterbewegung oscillations~\cite{schrodinger1930} at approximately four percent of the speed of light. The crucial insight lies in recognizing that the anomaly manifests differently depending on the observer's reference frame: while an observer co-moving with the internal motion measures a $g$-factor of exactly 2, an observer in the laboratory frame perceives Lorentz contraction effects that appear as the experimentally observed anomalous magnetic moment.

This observer-dependent interpretation fundamentally reframes the anomalous magnetic moment from an interaction-based quantum correction to a relativistic geometric effect arising from internal motion. The mathematical relationship between the measured anomaly and Lorentz contraction ratio establishes a direct connection between quantum magnetic phenomena and special relativity without requiring external field interactions or perturbative quantum field theory calculations. When extended to include general relativistic effects through geodetic precession~\cite{thomas1926,thomas1927}, this framework provides specific predictions for critical radii where different leptons maintain stable internal oscillations, offering geometric explanations for particle decay processes and generational hierarchy in the lepton family.

The implications extend beyond magnetic moments to encompass fundamental questions about the nature of quantum corrections and their relationship to relativistic geometry. If quantum phenomena traditionally attributed to virtual particle interactions can instead arise from observable relativistic effects of internal structure, this suggests a broader program for reinterpreting quantum mechanics through geometric principles. The present work develops this geometric framework systematically, demonstrating how observer-dependent effects can account for quantum magnetic anomalies while providing testable predictions for experimental verification.

\clearpage

% ----------------------------------------
% Table 1: Comparison of ZB-Based Models for Anomalous Magnetic Moment
% ----------------------------------------
\begin{table*}[t!]
\caption{Comparison of Zitterbewegung-Based Models for Anomalous Magnetic Moment}
\label{tab:zb_comparison}
\renewcommand{\arraystretch}{1.4}
\begin{tabular}{p{4.2cm}p{4.0cm}p{4.0cm}p{4.0cm}}
\toprule
\textbf{Aspect} & \textbf{Hestenes (2010)} & \textbf{Kozlov \& Nemchenko (2012)} & \textbf{0-Sphere Model (2018 -- 2025)} \\
\midrule
Anomalous magnetic moment treatment & Derives $g = 2$ only; $a_e$ not explicitly addressed & Semi-classical approximation for $a_e \approx 0.0011614$ & Exact relation $a_e = \gamma - 1$ with predictive capability \\
\addlinespace[0.3cm]
Zitterbewegung physical interpretation & Helical motion at light speed arising from Dirac algebra & Classical circular motion with approximate relativistic corrections & Subluminal thermal kernel oscillations at $v \sim 0.04c$ between spatially separated regions \\
\addlinespace[0.3cm]
Observer dependence & Not explicitly considered; mathematical feature of Dirac formalism & Limited treatment; focuses on fixed frame corrections & Central principle: anomaly vanishes in co-moving frame, appears in lab frame \\
\addlinespace[0.3cm]
Geometric structure & Sophisticated geometric algebra representation of spinor dynamics & Simple classical radius-velocity model without detailed structure & Lorentz contraction-based geometry with dual thermal kernels and photon sphere \\
\addlinespace[0.3cm]
Theoretical foundation & Reinterpretation of Dirac equation through geometric algebra & Semi-classical bridge between quantum and classical mechanics & Special relativistic observer transformations applied to internal motion \\
\addlinespace[0.3cm]
Experimental predictions & Visualization of ZB as real phenomenon; no specific velocity predictions & Specific orbital parameters for classical model & Direct prediction: internal velocity $v \sim 0.04c$, frequency $\sim 10^{18}$ Hz \\
\addlinespace[0.3cm]
External field requirements & Maintains standard QED assumption of field-dependent measurements & Follows conventional measurement approaches & Proposes field-free anomaly observable through internal motion detection \\
\addlinespace[0.3cm]
Mathematical approach & Advanced geometric algebra and multivector calculus & Classical mechanics with relativistic corrections & Direct application of Lorentz transformations: $\gamma = 1 + a_e$ \\
\addlinespace[0.3cm]
Scope of application & General framework for understanding electron structure & Specific to anomalous magnetic moment calculation & Anomalous magnetic moments plus broader geometric quantum phenomena \\
\bottomrule
\end{tabular}
\end{table*}


% ========================================
% Section 2: Motivation
% ========================================
\section{Motivation}

% ----------------------------------------
% Subsection 2.1: Fundamental Departure from QED
% ----------------------------------------
\subsection{Fundamental Departure from QED: Observer-Dependent Anomaly in Free Electrons}

The conventional quantum electrodynamics (QED) interpretation of the anomalous magnetic moment fundamentally assumes the necessity of external magnetic field interactions. In the standard QED framework, the total Hamiltonian is decomposed as 
\begin{equation}
H_{\text{total}} = H_0 + H_1,
\end{equation}
where $H_0$ represents the free particle Hamiltonian derived from the Dirac equation, and $H_1$ constitutes the interaction term with external electromagnetic fields. The anomalous magnetic moment emerges exclusively through perturbative calculations of $H_1$, implying that without external field interactions ($\mathbf{B} = \mathbf{0}$), no anomaly should manifest in the magnetic moment of a truly free electron. This interpretation suggests that the deviation of the $g$-factor from the Dirac value of 2 is purely a consequence of field-particle interactions, requiring the presence of $|\mathbf{B}| > 0$ for any observable anomaly to occur.

Our theoretical framework presents a fundamental challenge to this interpretation by proposing that the anomalous magnetic moment arises as an intrinsic geometric effect observable in free electrons ($\mathbf{B} = \mathbf{0}$) through relativistic observer transformations. The key insight lies in recognizing that the anomaly is not an interaction-dependent phenomenon but rather emerges from the relativistic relationship between different inertial reference frames observing the electron's internal structure. In the 0-Sphere model, the electron comprises two thermal kernels exchanging energy through a photon sphere exhibiting Zitterbewegung oscillations~\cite{schrodinger1930} at approximately $v \sim 0.04c$. The crucial distinction emerges when we consider different observational perspectives: an observer co-moving with the photon sphere's energy centroid experiences no relativistic distortion and measures a $g$-factor of exactly 2, consistent with the Dirac prediction. However, an observer in the laboratory frame, stationary relative to the oscillating internal structure, perceives Lorentz contraction effects that manifest as the experimentally observed anomalous magnetic moment.

This observer-dependent interpretation fundamentally reframes the anomalous magnetic moment from an interaction-based quantum correction to a geometric consequence of internal relativistic motion. The mathematical relationship 
\begin{equation}
    \frac{L}{L_0} = \frac{1}{1 + a_e},
    \label{eq:length_contraction}
\end{equation}
establishes a direct connection between the measured anomaly and the Lorentz contraction ratio, where $L_0$ represents the proper length in the co-moving frame and $L$ denotes the contracted length observed in the laboratory frame. Note that this relationship can be equivalently expressed as 
\begin{equation}
\gamma = 1 + a_e,
\label{eq:gamma_anomaly_relation}
\end{equation}
where $\gamma$ is the Lorentz factor, providing a more direct connection to relativistic transformations. When $a_e \to 0$, the internal velocity approaches zero ($v \to 0$), eliminating both the Thomas precession term and the observable anomaly, naturally recovering the Dirac value. Conversely, finite values of $a_e$ correspond to measurable internal velocities that produce relativistic effects observable only from external reference frames. This interpretation provides a first-principles explanation for why the anomalous magnetic moment appears at all, bridging the conceptual gap between quantum mechanical predictions and relativistic geometry without requiring external field interactions or perturbative quantum field theory calculations.

% ----------------------------------------
% Subsection 2.2: Comparison with Other Zitterbewegung-Based Models
% ----------------------------------------
\subsection{Comparison with Other Zitterbewegung-Based Models}

While the present framework introduces a novel observer-dependent geometric interpretation of the anomalous magnetic moment grounded in internal Zitterbewegung dynamics, it is not the first attempt to associate the $g$-factor anomaly with internal electron motion. Prior models, most notably by Hestenes~\cite{hestenes1990,hestenes2010}, have reinterpreted the Dirac equation's internal Zitterbewegung as a manifestation of helical motion in spacetime. In Hestenes' approach, the spin and magnetic moment arise from the phase and orientation of a lightlike helical path, eliminating the need for external field interactions. However, his formulation largely treats the internal motion as a mathematical feature of the Dirac algebra, without offering a specific mechanism for how relativistic transformations modify the measured $g$-factor between different observers.

Another class of semi-classical approaches, such as those proposed by Kozlov and Nemchenko~\cite{kozlov2012}, derive the anomalous magnetic moment by modeling Zitterbewegung as classical circular motion with relativistic corrections. Their results reproduce the leading-order anomaly, specifically the Schwinger term~\cite{schwinger1948} $a_e \approx \alpha/(2\pi) \approx 0.001161$, by assuming an effective radius and frequency consistent with internal oscillations, but without incorporating observer-dependent contraction or frame-specific symmetry. Similar efforts by Barut and coworkers~\cite{barut1981,barut1984} have explored classical models of electron structure that incorporate Zitterbewegung motion, though these approaches generally maintain fixed internal trajectories rather than emphasizing relativistic observer effects.

In contrast, the 0-Sphere model developed in this work emphasizes the observer-dependent nature of the anomaly as a fundamental geometric consequence of relativistic length contraction applied to subluminal internal oscillations. Rather than postulating fixed internal trajectories or relying on Dirac algebra alone, \textbf{the model links the measured deviation from $g = 2$ to the Lorentz contraction perceived by external observers.} This direct connection leads to a quantifiable prediction for internal Zitterbewegung velocity, $v \sim 0.04c$, which in turn determines the anomaly via the relation $\gamma = 1 + a_e$.

A systematic comparison of these different Zitterbewegung-based approaches reveals fundamental distinctions in their treatment of the anomalous magnetic moment, as summarized in Table~\ref{tab:zb_comparison}. While Hestenes' geometric algebra formulation provides elegant mathematical machinery for visualizing electron internal structure, it stops short of addressing the quantitative origin of the $g$-factor deviation from the Dirac value. Kozlov and Nemchenko's semi-classical model successfully reproduces the Schwinger term $a_e \approx \alpha/(2\pi) \approx 0.001161$ through classical orbital mechanics with relativistic corrections, but lacks a principled explanation for why such specific orbital parameters (e.g., Compton wavelength-scale radius) should arise or how they relate to observer-dependent relativistic principles.

The present framework distinguishes itself through its emphasis on observer-dependent geometric effects as the fundamental source of the anomaly. Unlike previous ZB-based models that treat the internal motion as a fixed physical feature leading to modified magnetic properties, our approach recognizes that the same internal dynamics yield different observational outcomes depending on the reference frame of measurement. This observer-dependent interpretation not only provides a natural explanation for the existence of the anomaly but also yields specific, testable predictions for the internal velocity scale ($v \sim 0.04c$) that directly connects relativistic kinematics to quantum magnetic phenomena.

Furthermore, while earlier models maintain the conventional assumption that external magnetic fields are necessary for anomaly detection, the geometric framework developed here suggests that \textbf{the anomalous magnetic moment represents an intrinsic property of free electrons observable through appropriate measurement of their internal dynamics.} This conceptual shift from interaction-dependent to geometry-dependent anomalies opens new experimental pathways for testing fundamental aspects of electron structure without relying on traditional field-based measurement techniques. Recent experimental advances in quantum simulation~\cite{gerritsma2010} and direct observation of Zitterbewegung in various physical systems~\cite{leblanc2013, schliemann2005, zawadzki2011, zhang2023} provide encouraging evidence that such internal motion phenomena may indeed be observable under appropriate conditions.

This formulation allows not only for a conceptually unified explanation of the anomaly but also offers a concrete path toward experimental validation through direct velocity measurement or high-frequency oscillation detection, distinguishing it from earlier ZB-based interpretations.

% ========================================
% Section 3: Observer-Dependent Anomalous Magnetic Moment
% ========================================
\section{Observer-Dependent Anomalous Magnetic Moment}

% ----------------------------------------
% Subsection 3.1: Internal ZB Dynamics and the Geometric Origin of Spin Anomaly
% ----------------------------------------
\subsection{Internal Zitterbewegung Dynamics and the Geometric Origin of Spin Anomaly}

The revelation that anomalous magnetic moments can arise from internal motion at relativistic frequencies, rather than external interactions, necessitates a comprehensive reexamination of electron spin dynamics. Traditional interpretations of Zitterbewegung~\cite{schrodinger1930,hestenes1990,recami1995} describe it as light-speed oscillatory motion arising from matter-antimatter interference, presenting theoretical challenges due to its superluminal character. Our model reinterprets Zitterbewegung as subluminal oscillation between two spatially separated thermal kernels within a single electron, occurring at frequencies on the order of $10^{18}$ Hz with average velocities around $0.04c$. This reinterpretation eliminates the need for electron-positron pair creation and instead describes the internal energy exchange process as $e^-_{\text{kernelA}} \to \gamma^*_{K.E.} \to e^-_{\text{kernelB}}$~\cite{hanamura2309.0047}, where the photon sphere converts thermal potential energy from one kernel into kinetic energy subsequently absorbed by the second kernel.

The geometric foundation of this model rests on the recognition that non-circular internal motion can generate angular momentum through appropriate phase relationships, challenging the classical intuition that intrinsic spin must originate from rigid-body rotation. The Thomas precession formula~\cite{thomas1926,thomas1927}
\begin{equation}
    \boldsymbol{\Omega} = \frac{1}{2c^2}[\mathbf{a} \times \mathbf{v}],
    \label{eq:Thomas1}
\end{equation}
applies generally to any system exhibiting internal acceleration and velocity components, including reciprocating motions. When $\boldsymbol{v}$ is much smaller than the speed of light $c$, the resulting precession frequency $\boldsymbol{\Omega}$ is negligible due to the $c^2$ term in the denominator of Eq.~\eqref{eq:Thomas1}.

This explains why classical systems such as pendulums, spring oscillators, or any low-velocity mechanical vibrators observable in daily life show no measurable spin-related relativistic effects. This phenomenon parallels other scale-dependent physical effects: just as de Broglie wavelengths become negligible for macroscopic masses (eliminating quantum effects for human-scale objects), and Einstein's field equations reduce to Newtonian mechanics in weak gravitational fields~\cite{jackson1999}, the Thomas precession becomes significant only when internal velocities approach relativistic scales. The crucial premise of the 0-Sphere model emerges from this reductionist principle: when Zitterbewegung micro-oscillation velocities approach the speed of light, angular velocity generation becomes non-negligible, as dictated by the mathematical structure of Eq.~\eqref{eq:Thomas1}.

However, based on the relation given in Eq.~\eqref{eq:length_contraction}, a nonzero anomalous magnetic moment $a_e$ implies a finite contraction ratio and hence a finite velocity $v \sim 0.04c$. At this velocity scale, relativistic corrections such as Thomas precession become non-negligible, naturally giving rise to angular momentum and the observed $g$-factor anomaly.

When applied to harmonic oscillation with $v(t) = v_0\cos(\omega t)$ and $a(t) = -v_0\omega\sin(\omega t)$, the cross product yields angular velocity proportional to $\sin(2\omega t)$, revealing the emergence of double-frequency precession from single-frequency oscillation. This mathematical result provides a natural explanation for why spin angular momentum quantized to $\hbar/2$ generates magnetic moments equivalent to those produced by full orbital angular momentum $\hbar$, resolving the long-standing puzzle of spin's enhanced magnetic efficiency first noted in early spin studies~\cite{uhlenbeck1925, uhlenbeck1926}.

The observer-dependent nature of this phenomenon becomes apparent when examining the system from different reference frames. In the co-moving frame of the internal oscillation, where an observer travels with the Zitterbewegung motion, the internal structure appears symmetric and non-relativistic, yielding a $g$-factor of precisely 2 as predicted by the Dirac equation. The internal motion exhibits no net angular momentum or magnetic anomaly from this perspective, since the observer experiences no relative motion with respect to the oscillating system.

However, from the laboratory frame, where the observer remains stationary while the internal structure oscillates at relativistic speeds, Lorentz contraction distorts the effective geometry of the motion. This distortion manifests as a measurable deviation in the magnetic moment, with the magnitude of the anomaly directly proportional to the degree of relativistic contraction experienced by the internal oscillation. The relationship $\gamma = 1 + a_e$ quantitatively captures this geometric transformation, establishing a direct bridge between special relativity and quantum magnetic phenomena without invoking virtual particle interactions or external field perturbations.

% ----------------------------------------
% Subsection 3.2: Implications for QM and Relativistic Consistency
% ----------------------------------------
\subsection{Implications for Quantum Mechanics and Relativistic Consistency}

The interpretation of anomalous magnetic moments as observer-dependent geometric effects carries profound implications for our understanding of the interface between quantum mechanics and special relativity. This framework suggests that certain quantum mechanical phenomena traditionally attributed to field interactions may actually represent relativistic artifacts arising from the frame-dependent observation of internal structure. According to this interpretation, the anomaly disappears in the co-moving frame and emerges only in the laboratory frame, implying that \textbf{the $g$-factor should be regarded as a relational quantity rather than an intrinsic property}---fundamentally dependent on the observer's state of motion relative to the particle's internal dynamics. This perspective aligns with the principle of relativity while providing an intuitive explanation for the existence of quantum corrections that have historically required sophisticated perturbative calculations to understand.

The consistency of this interpretation with the running of the fine structure constant $\alpha$ provides additional theoretical support for the geometric approach to quantum phenomena. As energy scales increase, $\alpha$ evolves from approximately $1/137$ at low energies to $\sim 1/128$ at higher energies, suggesting a fundamental connection between electromagnetic coupling strength and the energy scale of internal motion. In our framework, higher energy corresponds to faster Zitterbewegung velocities, which in turn produce stronger relativistic effects and larger observable anomalies. This natural correlation between coupling strength evolution and internal motion dynamics suggests that the geometric interpretation may capture deeper aspects of quantum field behavior than previously recognized. Furthermore, the framework's prediction that electrons maintain stable Zitterbewegung oscillations while heavier leptons reach critical radii~\cite{hanamura2501.0130} where such motion becomes unsustainable provides a geometric explanation for the observed pattern of lepton stability and decay processes.

The absence of external magnetic field requirements in this model fundamentally challenges the conventional understanding that anomalous magnetic moments arise solely from particle-field interactions. By demonstrating that such anomalies can emerge from pure geometric effects in free particles, this framework opens new avenues for understanding quantum phenomena as manifestations of relativistic internal structure rather than products of field theory interactions. This shift in perspective suggests that many quantum corrections traditionally calculated through perturbative methods may have underlying geometric origins that become apparent when viewed from the appropriate relativistic framework. The implications extend beyond magnetic moments to potentially encompass other quantum phenomena that exhibit observer-dependent characteristics, suggesting a broader program for reinterpreting quantum mechanics through the lens of relativistic geometry and internal particle dynamics.

Despite the theoretical elegance of interpreting anomalous magnetic moments as observer-dependent geometric effects in free electrons, their direct experimental verification remains fundamentally constrained by current measurement techniques. Specifically, the conventional reliance on external magnetic fields to induce and detect spin precession highlights a critical limitation: without a nonzero magnetic field ($\mathbf{B} = \mathbf{0}$), standard QED-based apparatuses lack the operational means to measure deviations in the $g$-factor. This constraint renders the experimental observation of free-electron anomalies intrinsically challenging, if not currently infeasible. As a viable alternative, the focus must shift toward the direct measurement of the internal Zitterbewegung velocity, which the geometric model predicts to be approximately $v \sim 0.04c$. Since this velocity serves as the underlying source of the relativistic contraction effects leading to the anomaly, its precise experimental determination would provide a critical test of the geometric framework. Advancing techniques capable of probing such rapid substructure dynamics---potentially through high-frequency Compton scattering or ultrafast temporal interference---thus becomes an urgent priority for validating the model in the absence of external field interactions.

This theoretical framework suggests that many quantum phenomena traditionally attributed to field interactions may represent geometric consequences of internal structure observed from inappropriate reference frames. The ability to construct a theory where free electrons exhibit anomalous magnetic moments without external field requirements represents a significant conceptual advance, potentially opening new avenues for understanding quantum corrections as manifestations of relativistic geometry rather than products of perturbative field theory calculations. The observer-dependent nature of magnetic anomalies indicates that quantum mechanical properties should be treated as relational quantities fundamentally dependent on the observer's state of motion relative to internal particle dynamics.

% ----------------------------------------
% Subsection 3.3: Conditions for Zero Zitterbewegung Velocity
% ----------------------------------------
\subsection{Conditions for Zero Zitterbewegung Velocity}

The theoretical framework developed here predicts two distinct physical conditions under which the Zitterbewegung velocity approaches zero, each arising from fundamentally different relativistic considerations. The first condition emerges within the special relativistic framework and applies universally to all leptons, while the second condition involves general relativistic effects and applies specifically to unstable particles undergoing decay processes.

The primary condition for vanishing Zitterbewegung velocity occurs in the co-moving reference frame, where an observer travels with the energy centroid of the photon sphere mediating energy exchange between thermal kernels. In this reference frame, the internal oscillatory motion appears stationary, eliminating relativistic distortions and yielding a $g$-factor of exactly 2 consistent with Dirac theory. This condition represents a fundamental consequence of special relativistic observer transformations and applies equally to electrons, muons, and tau particles, as each lepton maintains internal structure characterized by dual thermal kernels and photon sphere dynamics. The universal applicability of this condition reflects the underlying geometric nature of the anomalous magnetic moment, which emerges solely from the relativistic relationship between different inertial reference frames observing the same internal dynamics.

The second condition for zero Zitterbewegung velocity emerges from general relativistic considerations involving geodetic precession effects that determine particle stability thresholds. According to our previous work~\cite{hanamura2501.0130} extending the framework to incorporate Thomas precession and gravitational effects, heavier leptons reach critical radii where their internal oscillations become unsustainable, triggering decay processes into lighter particles. At the moment of decay, when a muon or tau particle transitions to its decay products, the Zitterbewegung velocity of the decaying particle approaches zero as its internal structure reaches the critical radius limit. This condition differs fundamentally from the co-moving frame scenario, as it represents a physical threshold determined by the interplay between internal motion dynamics and relativistic precession effects rather than a choice of reference frame. The decay condition applies specifically to muons and tau particles, which possess sufficient mass to reach critical radii, while electrons remain stable below such thresholds and do not undergo further decay into lighter leptons.

These two distinct conditions for vanishing Zitterbewegung velocity illuminate different aspects of the geometric interpretation of particle physics. The co-moving frame condition demonstrates the observer-dependent nature of quantum magnetic properties within special relativity, while the decay threshold condition reveals how general relativistic effects can determine particle stability through internal structure limitations. The coexistence of these conditions within a unified geometric framework suggests that both quantum mechanical and particle physics phenomena may have deeper connections to relativistic geometry than previously recognized, providing a foundation for understanding diverse physical processes through consistent geometric principles.


% ========================================
% Section 4: Conclusion
% ========================================
\section{Conclusion}

The analysis presented in this work establishes that anomalous magnetic moments can arise as purely geometric effects in free electrons, independent of external magnetic field interactions. This conclusion fundamentally challenges the conventional quantum electrodynamics interpretation that requires external field perturbations for any observable magnetic anomaly to manifest. The key to this reinterpretation lies in the 0-Sphere model's internal structure, which enables simultaneous consideration of both co-moving and external observer perspectives within a single electron description through its harmonic oscillator dynamics mediated by energy exchange between thermal kernels and a photon sphere.

The key insight lies in the dual perspective framework enabled by the 0-Sphere model's internal structure. An observer co-moving with the photon sphere's energy centroid experiences no relativistic distortion and measures a $g$-factor of exactly 2, consistent with the Dirac prediction for a truly free particle. However, an external observer in the laboratory frame perceives the accelerating and decelerating motion of the photon sphere, experiencing Lorentz contraction effects that manifest as the experimentally observed anomalous magnetic moment. This observer-dependent phenomenon occurs within the framework of special relativity, as the distinction between co-moving and external reference frames requires no general relativistic considerations for its basic formulation.

The mathematical relationship $\gamma = 1 + a_e$ establishes a direct algebraic connection between measured anomalies and relativistic geometry, demonstrating that quantum corrections can emerge from pure geometric effects observable in free particles. This formulation eliminates the conceptual requirement that anomalous magnetic moments arise only through external field interactions, instead showing that such anomalies represent natural consequences of internal structure observed from appropriate reference frames. The framework predicts specific internal velocities around $v \sim 0.04c$ for electron Zitterbewegung, providing concrete targets for experimental verification.

The 0-Sphere model's incorporation of harmonic oscillator dynamics within electron structure enables this dual-perspective interpretation by providing physical mechanisms for both co-moving and external observations within a single particle description. The photon sphere's simple harmonic motion creates the necessary internal reference frame for co-moving observations while simultaneously providing the accelerating motion observable from external frames. Although the complete theoretical framework incorporates general relativistic effects through geodetic precession~\cite{thomas1926,thomas1927} to determine critical radii for different leptons, the fundamental principle of observer-dependent anomalous magnetic moments operates entirely within special relativistic considerations.

This theoretical framework suggests that many quantum phenomena traditionally attributed to field interactions may represent geometric consequences of internal structure observed from external reference frames. The ability to construct a theory where free electrons exhibit anomalous magnetic moments without external field requirements represents a significant conceptual advance, potentially opening new avenues for understanding quantum corrections as manifestations of relativistic geometry rather than products of perturbative field theory calculations. The observer-dependent nature of magnetic anomalies indicates that quantum mechanical properties should be treated as relational quantities fundamentally dependent on the observer's state of motion relative to internal particle dynamics.

% ========================================
% Bibliography
% ========================================
\begin{thebibliography}{99}

\bibitem{fan2023}
X.~Fan, T.~G.~Myers, B.~A.~D.~Sukra, and G.~Gabrielse, 
``Measurement of the Electron Magnetic Moment,'' 
\textit{Phys. Rev. Lett.} \textbf{130}, 071801 (2023). 
\href{https://doi.org/10.1103/PhysRevLett.130.071801}{https://doi.org/10.1103/PhysRevLett.130.071801}

\bibitem{hanneke2008}
D.~Hanneke, S.~Fogwell, and G.~Gabrielse, 
``New Measurement of the Electron Magnetic Moment and the Fine Structure Constant,'' 
\textit{Phys. Rev. Lett.} \textbf{100}, 120801 (2008). 
\href{https://doi.org/10.1103/PhysRevLett.100.120801}{https://doi.org/10.1103/PhysRevLett.100.120801}

\bibitem{gabrielse2006}
G.~Gabrielse, D.~Hanneke, T.~Kinoshita, M.~Nio, and B.~Odom, 
``New Determination of the Fine Structure Constant from the Electron g Value and QED,'' 
\textit{Phys. Rev. Lett.} \textbf{97}, 030802 (2006). 
\href{https://doi.org/10.1103/PhysRevLett.97.030802}{https://doi.org/10.1103/PhysRevLett.97.030802}

\bibitem{schwinger1948}
J.~Schwinger, 
``On Quantum-Electrodynamics and the Magnetic Moment of the Electron,'' 
\textit{Phys. Rev.} \textbf{73}, 416--417 (1948). 
\href{https://doi.org/10.1103/PhysRev.73.416}{https://doi.org/10.1103/PhysRev.73.416}

\bibitem{feynman1948}
R.~P.~Feynman, 
``Space-Time Approach to Quantum Electrodynamics,'' 
\textit{Phys. Rev.} \textbf{76}, 769--789 (1949). 
\href{https://doi.org/10.1103/PhysRev.76.769}{https://doi.org/10.1103/PhysRev.76.769}

\bibitem{feynman1949}
R.~P.~Feynman, 
``The Theory of Positrons,'' 
\textit{Phys. Rev.} \textbf{76}, 749--759 (1949). 
\href{https://doi.org/10.1103/PhysRev.76.749}{https://doi.org/10.1103/PhysRev.76.749}

\bibitem{tomonaga1946}
S.~Tomonaga, 
``On a relativistically invariant formulation of the quantum theory of wave fields,'' 
\textit{Prog. Theor. Phys.} \textbf{1}, 27--42 (1946). 
\href{https://doi.org/10.1143/PTP.1.27}{https://doi.org/10.1143/PTP.1.27}

\bibitem{aoyama2015}
T.~Aoyama, M.~Hayakawa, T.~Kinoshita, and M.~Nio, 
``Tenth-Order Electron Anomalous Magnetic Moment,'' 
\textit{Phys. Rev. D} \textbf{91}, 033006 (2015). 
\href{https://doi.org/10.1103/PhysRevD.91.033006}{https://doi.org/10.1103/PhysRevD.91.033006}

\bibitem{dirac1928}
P.~A.~M.~Dirac, 
``The quantum theory of the electron,'' 
\textit{Proc. R. Soc. Lond. A} \textbf{117}, 610--624 (1928). 
\href{https://doi.org/10.1098/rspa.1928.0023}{https://doi.org/10.1098/rspa.1928.0023}

\bibitem{hanamura2018}
S.~Hanamura, A Model of an Electron Including Two Perfect Black Bodies, Zenodo (2018), \href{https://doi.org/10.5281/zenodo.16759284}{doi:10.5281/zenodo.16759284}. [Originally: \href{https://vixra.org/abs/1811.0312}{viXra:1811.0312}]

\bibitem{schrodinger1930}
E.~Schr{\"o}dinger, 
``{\"U}ber die kr{\"a}ftefreie Bewegung in der relativistischen Quantenmechanik,'' 
\textit{Sitzungsber. Preuss. Akad. Wiss. Phys. Math. Kl.} \textbf{24}, 418--428 (1930).

\bibitem{thomas1926}
L.~H.~Thomas, 
``The Motion of the Spinning Electron,'' 
\textit{Nature} \textbf{117}, 514 (1926); 
``The kinematics of an electron with an axis,''
\textit{Phil. Mag.} \textbf{3}, 1 (1927). 
\href{https://doi.org/10.1080/14786440408564502}{https://doi.org/10.1080/14786440408564502}

\bibitem{thomas1927}
L.~H.~Thomas, 
``The kinematics of an electron with an axis,'' 
\textit{Phil. Mag.} \textbf{3}, 1--22 (1927). 
\href{https://doi.org/10.1080/14786440108564170}{https://doi.org/10.1080/14786440108564170}

\bibitem{hestenes1990}
D.~Hestenes, 
``The zitterbewegung interpretation of quantum mechanics,'' 
\textit{Found. Phys.} \textbf{20}, 1213--1232 (1990). 
\href{https://doi.org/10.1007/BF01889466}{https://doi.org/10.1007/BF01889466}

\bibitem{hestenes2010}
D.~Hestenes, 
``Zitterbewegung in Quantum Mechanics,'' 
\textit{Found. Phys.} \textbf{40}, 1 (2010). 
\href{https://doi.org/10.1007/s10701-009-9360-3}{https://doi.org/10.1007/s10701-009-9360-3}

\bibitem{kozlov2012}
V.~V.~Kozlov and E.~E.~Nemchenko, 
``Semi-classical explanation of the anomalous magnetic moment of the electron,'' 
\textit{J. Phys.: Conf. Ser.} \textbf{343}, 012123 (2012).
\href{https://doi.org/10.1088/1742-6596/343/1/012123}{https://doi.org/10.1088/1742-6596/343/1/012123}

\bibitem{barut1981}
A.~O.~Barut and A.~J.~Bracken, 
``Zitterbewegung and the internal geometry of the electron,'' 
\textit{Phys. Rev. D} \textbf{23}, 2454--2463 (1981). 
\href{https://doi.org/10.1103/PhysRevD.23.2454}{https://doi.org/10.1103/PhysRevD.23.2454}

\bibitem{barut1984}
A.~O.~Barut and N.~Zanghi, 
``Classical model of the Dirac electron,'' 
\textit{Phys. Rev. Lett.} \textbf{52}, 2009--2012 (1984).
\href{https://doi.org/10.1103/PhysRevLett.52.2009}{https://doi.org/10.1103/PhysRevLett.52.2009}

\bibitem{gerritsma2010}
R.~Gerritsma, G.~Kirchmair, F.~Z{\"a}hringer, E.~Solano, R.~Blatt, and C.~F.~Roos, 
``Quantum simulation of the Dirac equation,'' 
\textit{Nature} \textbf{463}, 68--71 (2010). 
\href{https://doi.org/10.1038/nature08688}{https://doi.org/10.1038/nature08688}

\bibitem{leblanc2013}
L.~J.~LeBlanc, M.~C.~Beeler, K.~Jim{\'e}nez-Garc{\'i}a, A.~R.~Perry, S.~Sugawa, R.~A.~Williams, and I.~B.~Spielman, 
``Direct observation of Zitterbewegung in a Bose–Einstein condensate,'' 
\textit{New J. Phys.} \textbf{15}, 073011 (2013). 
\href{https://doi.org/10.1088/1367-2630/15/7/073011}{https://doi.org/10.1088/1367-2630/15/7/073011}

\bibitem{schliemann2005}
J.~Schliemann, D.~Loss, and R.~M.~Westervelt, 
``Zitterbewegung of electronic wave packets in III-V zinc-blende semiconductor quantum wells,'' 
\textit{Phys. Rev. Lett.} \textbf{94}, 206801 (2005). 
\href{https://doi.org/10.1103/PhysRevLett.94.206801}{https://doi.org/10.1103/PhysRevLett.94.206801}

\bibitem{zawadzki2011}
W.~Zawadzki and T.~M.~Rusin, 
``Zitterbewegung (trembling motion) of electrons in semiconductors: a review,'' 
\textit{J. Phys.: Condens. Matter} \textbf{23}, 143201 (2011). 
\href{https://doi.org/10.1088/0953-8984/23/14/143201}{https://doi.org/10.1088/0953-8984/23/14/143201}

\bibitem{zhang2023}
Z.~Q.~Zhang, Y.~V.~Kartashov, F.~Ye, Y.~Y.~Lin, A.~Szameit, and Z.~Chen, 
``Observation of Zitterbewegung in photonic microcavities,'' 
\textit{Light: Sci. Appl.} \textbf{12}, 289 (2023). 
\href{https://doi.org/10.1038/s41377-023-01177-z}{https://doi.org/10.1038/s41377-023-01177-z}

\bibitem{recami1995}
E.~Recami and G.~Salesi, 
``Field Theory of the Electron, Spin and Zitterbewegung,'' 
arXiv:hep-th/9508168 (1995). 
\href{https://arxiv.org/abs/hep-th/9508168}{https://arxiv.org/abs/hep-th/9508168}

\bibitem{hanamura2309.0047} 
S.~Hanamura, Redefining Electron Spin and Anomalous Magnetic Moment Through Harmonic Oscillation and Lorentz Contraction, Zenodo (2023), \href{https://doi.org/10.5281/zenodo.17764997}{doi:10.5281/zenodo.17764997}. [Originally: \href{https://vixra.org/abs/2309.0047}{viXra:2309.0047}]

\bibitem{jackson1999}
J.~D.~Jackson, 
\textit{Classical Electrodynamics}, 3rd ed., 
Wiley, New York (1999).

\bibitem{uhlenbeck1925}
G.~E.~Uhlenbeck and S.~A.~Goudsmit, 
``Spinning Electrons and the Structure of Spectra,'' 
\textit{Naturwissenschaften} \textbf{13}, 953 (1925).
\href{https://doi.org/10.1007/BF01558878}{https://doi.org/10.1007/BF01558878}

\bibitem{uhlenbeck1926}
G.~E.~Uhlenbeck and S.~A.~Goudsmit, 
``Spinning Electrons and the Structure of Spectra,'' 
\textit{Nature} \textbf{117}, 264 (1926).
\href{https://doi.org/10.1038/117264a0}{https://doi.org/10.1038/117264a0}

\bibitem{hanamura2501.0130} 
S.~Hanamura, Bridging Quantum Mechanics and General Relativity: A First-Principles Approach to Anomalous Magnetic Moments and Geodetic Precession, Zenodo (2025), \href{https://doi.org/10.5281/zenodo.17765097}{doi:10.5281/zenodo.17765097}. [Originally: \href{https://vixra.org/abs/2501.0130}{viXra:2501.0130}]

\end{thebibliography}

\end{document}